\section{The 49 versions}
\label{app:ledger}

\begin{table*}[t]
\centering
\small
\begin{tabularx}{\textwidth}{P{0.12\linewidth}P{0.22\linewidth}YY}
\toprule
Versions & Phase & Kept & Ruled out \\
\midrule
v001--v005 & Lie-group basics: $SO(3)$ benchmarks, corrected right-action CF4/RKMK4, Yoshida
composition, two-stage Gauss--Lie & Gauss collocation as the high-order family; right-action
conventions & Yoshida composition (negative substeps unsuitable for friction) \\
v006--v013 & Constrained DAE closure: reduced fixed-pivot DAE, absolute-coordinate Gauss DAE,
endpoint constraints inside Newton, friction, AD Jacobian, quaternion transport, three-stage Gauss
& absolute coordinates; endpoint constraints in the solve; JAX automatic differentiation;
quaternions; six stages & position-only closure; finite-difference Jacobians \\
v014--v022 & Friction regimes and baselines: smoothness sweep, adaptive Gauss6/4, trapezoidal,
BDF2 and Lobatto baselines, multiplier-dependent Brown--McPhee friction & Brown--McPhee law through
the multipliers; the smooth/sharp split as a stated regime & trapezoidal, BDF2 and Lobatto endpoint
methods as accuracy candidates; adaptivity as the main result \\
v023--v029 & Absolute-coordinate revolute DAE: full DAE, Lobatto node swap, projection repair,
stage velocity rows, stage acceleration rows, off-axis torque, two joints & \emph{stage-level
velocity and acceleration constraint rows} (``FullVA''), the decisive move & node swap; post-step
projection as a fix \\
v030--v045 & Sparse solver scaling: Jacobian sparsity, sparse Newton, matrix-free Krylov, lagging,
coloured AD, symbolic patterns, caches, larger topologies, prismatic joints & correctness of
sparse patterns; prismatic and skew-axis coverage & every speed claim: dense AD stayed faster \\
v046--v048 & Validation anchors: four ASME mechanisms, two-body cylindrical chain pipeline,
cross-paper benchmark harness & the chain as the main benchmark; the public suite as the external
one & external superiority claims (source-policy rows 0/40) \\
v049 & Experiments for the rewritten paper & E1--E7 & long horizons ($T=1$, $T=20$) on the chain \\
\bottomrule
\end{tabularx}
\caption{Phases of the exploration. ``Kept'' is what survived into the final method; ``ruled out''
is preserved as a versioned negative result.}
\label{tab:phases}
\end{table*}

Table~\ref{tab:phases} condenses the exploration into phases; Table~\ref{tab:ledger} gives one line per version: the method family tested and the role the
version plays in the record (its ``current role'' field in the machine-readable ledger).

\begin{table*}[t]
\centering
\scriptsize\setlength{\tabcolsep}{3pt}
\begin{tabular}[t]{@{}p{0.025\textwidth}p{0.19\textwidth}p{0.245\textwidth}@{}}
\toprule
v & Method family & Role in the record \\
\midrule
001 & $SO(3)$ kinematics and Euler top & benchmark infrastructure \\
002 & CF4 and RKMK4 (right action) & prescribed-orientation baseline, order four \\
003 & public rA code reproduction & reproducibility anchor for the public suite \\
004 & Yoshida-composed Lie midpoint & conservative baseline; negative substeps unsuitable for friction \\
005 & two-stage Gauss--Lie & smooth conservative mechanics, order four \\
006 & reduced fixed-pivot Gauss DAE & reduced constrained prototype \\
007 & absolute-coordinate Gauss DAE & index-3 consistency lesson \\
008 & endpoint-constrained Gauss DAE & projection-free absolute-coordinate prototype \\
009 & frictional endpoint Gauss4 & friction order-reduction evidence \\
010 & JAX AD Jacobian & nonlinear-solve backend \\
011 & $S^3$ transport operators & bridge to the quaternion formulation \\
012 & quaternion endpoint Gauss4 & verified quaternion residual \\
013 & Gauss6 endpoint collocation & first smooth sixth-order candidate \\
014 & Gauss4 vs Gauss6 friction sweep & smoothness decision rule \\
015 & adaptive embedded Gauss6/4 & near-nonsmooth regime tool \\
016 & reduced Lie-trapezoidal & ruled out as accuracy candidate \\
017 & reduced Lie-BDF2 & ruled out as accuracy candidate \\
018 & reduced Lie-Lobatto endpoint & TFE-style baseline ruled out \\
019 & multiplier-dependent Stribeck friction & AD-friendly friction through multipliers \\
020 & adaptive multiplier friction & sharp-friction regime tool \\
021 & Brown--McPhee multiplier friction & the friction law kept \\
022 & revolute Brown--McPhee Gauss6/4 & paper-like revolute friction evidence \\
023 & absolute-coordinate revolute Gauss DAE & full DAE; endpoint drift exposed \\
024 & direct Lobatto full revolute DAE & node swap ruled out \\
025 & projected absolute revolute DAE & projection repair: works, changes the method \\
\bottomrule
\end{tabular}\hfill
\begin{tabular}[t]{@{}p{0.025\textwidth}p{0.19\textwidth}p{0.245\textwidth}@{}}
\toprule
v & Method family & Role in the record \\
\midrule
026 & stage pivot-velocity rows & non-projection closure, velocity level \\
027 & stage velocity + acceleration rows & \textbf{FullVA}: the stage system kept \\
028 & off-axis FullVA revolute & axis-torque stress test \\
029 & double-revolute FullVA & generalizes to two joints \\
030 & Jacobian sparsity diagnosis & motivates sparse solvers \\
031 & sparse Newton & correct; not faster \\
032 & matrix-free Newton--Krylov & ruled out \\
033 & lagged sparse Newton & ruled out \\
034 & coloured-JVP sparse Jacobian & correct structured sparse AD \\
035 & batched coloured JVP & sparse-AD backend \\
036 & block-symbolic sparse pattern & production-style pattern \\
037 & JVP-pruned pattern & pattern discovery \\
038 & pattern cache reuse & cache policy \\
039 & triple-revolute scaling & transfer to larger topology \\
040 & generated block triple pattern & pattern acquisition \\
041 & triple-revolute cache refresh & cache policy \\
042 & skew-axis triple revolute & non-coplanar geometry \\
043 & skew prismatic pair & first non-revolute pair \\
044 & double-prismatic chain & interbody prismatic pair \\
045 & row-coloured VJP sparse AD & sparse-assembly diagnostic \\
046 & four ASME mechanisms & public-suite validation anchor \\
047 & two-body cylindrical chain pipeline & the main benchmark and the accepted implementation \\
048 & cross-paper same-test harness & external comparison scaffold; 0/40 source-policy rows \\
049 & paper experiments E1--E7 & the numerical section of the rewritten manuscript \\
\bottomrule
\end{tabular}
\caption{Version ledger (condensed from \texttt{validation/docs/version\_ledger.csv}).}
\label{tab:ledger}
\end{table*}

\section{Memory files}
\label{app:memory}

The eleven memory files at the end of the program, by their one-line descriptions: where the Lean
project lives and why it is kept off the synchronized drive; what the Lean development covers and
what it must not be claimed to close; the row-family mismatch between manuscript and code and its
resolution ($F(Z_G)=0$ exactly); the manuscript compaction of September~17 and the builder order;
the implemented predictor's zero angular guesses and their consequence for the Newton lemma; the
derived arXiv version and how it is rebuilt; the five-folder repository layout and its path
routing; the guarded external campaign of September~20 (approval sentence, 0/40 rows, why it
cannot close by execution); the paper rewrite of September~20--21 and its state; the benchmark's
regular-branch limit and the choice $T=0.5$; and the driven single-pendulum harness whose predictor
is the analytic drive. Each file has a ``why'' and a ``how to apply'' paragraph; the index is
loaded at the start of every session.

\section{Human messages of the final session}
\label{app:messages}

The \nHuman{} messages typed by the human (the other \nHarness{} user-role entries of the
transcript are harness-generated: task notifications, command echoes, compaction summaries and
image attachments; \nShell{} more are shell commands typed for the GitHub login), lightly
paraphrased from Chinese and grouped. \emph{Goals} (12): can Lean prove the convergence order;
install it; run the Lean check and report; continue and look for problems; tidy the Lean
development; make an arXiv version in a simple template; reorganize the folders; split the LaTeX
out of the ledger; make a clean five-folder layout keeping only the best manuscript; read it again
and find what else is wrong; write this paper; keep fixing. \emph{Approvals of the agent's own
proposals} (9): ``do as you said'' and variants, including the two approvals of the rewrite plan.
\emph{Authorizations} (3): any permission including network access; push to GitHub as you see fit;
run the external campaign. \emph{Decisions} (2): option A for the benchmark horizon, against the
agent's recommendation; target NAACL rather than a machine-learning conference for this paper.
\emph{Judgments} (4): the paper quality feels poor; the PDF does not look good, visualize the
pages; why does the top-left panel of Figure~7 look strange; this paper's quality is very poor,
the figures and the content are unclear. \emph{Questions} (6): what is this doing; what do you
think is appropriate; what does that mean; what is still missing; what is the status; what else is
there to do. \emph{Other} (3): an acknowledgment, ``continue'', and a report that the GitHub login
had failed.

\section{The incidents in full}
\label{app:incidents}

Table~\ref{tab:incidents_full} gives, for each incident of Table~\ref{tab:incidents}, the
wrong statement, its cause, the detector and the resolution.

\begin{table*}[t]
\centering
\scriptsize
\begin{tabularx}{\textwidth}{@{}P{0.02\linewidth}P{0.27\linewidth}P{0.25\linewidth}P{0.15\linewidth}Y@{}}
\toprule
\# & Wrong or misleading statement & Cause & Detected by & Resolution \\
\midrule
1 & The manuscript was sound but unreadable: hundreds of ``diagnostic'', ``non-claim'' and
``boundary'' sentences. & Every validator-pinned phrase had been written to satisfy the pin. & H
(``the paper quality feels poor'') & Full rewrite; ledger frozen; one number-checking validator. \\
2 & Orientation error converged at order 2.2 while position, velocity and angular velocity
converged at order 6.2. & Geodesic angle computed as $2\arccos(p\cdot q)$, whose resolution floor
is $3\times10^{-8}$ rad. & C (one of four quantities disagreed) & Chord formula
$2\arcsin(\lVert p\mp q\rVert/2)$; orientation order 6.28. \\
3 & Reference runs on the chain failed at $h=0.1/256$ over $T=1$; the plan said $T=1$ and
$T=20$. & The benchmark's frozen sliding direction degenerates at $t\approx0.6$,
a modelling limit, not a solver failure. & A (Newton residual $6\times10^{5}$) then H (chose to
keep the model and stop at $T=0.5$) & All chain experiments on $[0,0.5]$; limit stated in the
paper. \\
4 & Driven single pendulum: velocity error rose and fell erratically between $10^{-15}$ and
$4\times10^{-9}$ while position showed order 6. & The harness's Newton predictor is the analytic
solution and the stage Jacobian has condition number $10^{10}$--$10^{13}$; velocities carry
amplified roundoff. & H (``why does the top-left panel look strange?'') & Velocity not reported for
that mechanism; the panel labelled as a reconstruction check; mechanism explained. \\
5 & The four ASME mechanisms were planned as four convergence tests. & Three of the four are
kinematically driven; their motion is fixed by the drive. & A (errors at $10^{-14}$ for all $h$) &
Only the double pendulum is a dynamics test; the others are consistency rows. \\
6 & The public-code comparison at $T=3$ showed the public methods with $O(1)$ errors; suspicion of
a model mismatch. & The double pendulum amplifies perturbations over $T=3$. & C (added a $T=0.1$
alignment run: first-order decrease, same model) & Comparison kept with the alignment check
reported. \\
7 & ``The $h^7$-scaled stopping rule with $c_\eta=10^{4}$ covers the reported grids.'' & With the
fixed tolerance $10^{-11}$, the accepted residual over $h^7$ reaches $10^{6}$ at the finest
step. & A (purpose-built experiment E7) & Statement replaced by the measured ratios and a
discussion of the pessimistic bound. \\
8 & The single pendulum ``converges at order 6 with a floor'' was fitted over roundoff points,
giving 3.6. & Fit used a zero floor for the analytic reference. & C & Roundoff floor $10^{-15}$;
fit 6.0 over the resolved points. \\
9 & ``All results generated from a single implementation path.'' & Four harnesses assemble the
same stage system for four mechanisms. & C (self-review) & Reworded; assemblies listed. \\
10 & After the rewrite the frozen ledger reported stale line counts and lost markers. & Its
snapshot records pinned the sizes of scripts that had changed. & A (frozen validators) & Builder
sequence rerun; both frozen chains green. \\
\bottomrule
\end{tabularx}
\caption{The ten incidents of the final phase in full (Table~\ref{tab:incidents} of the main text is the condensed version). ``Detected by'' is the mechanism that first raised the
problem; A = artefact check (validator or purpose-built experiment), C = agent cross-check, H =
human judgment.}
\label{tab:incidents_full}
\end{table*}

\section{Artefacts and cost of the checks}
\label{app:cost}

\begin{table}[H]
\centering
\small
\begin{tabular}{@{}lr@{}}
\toprule
Versioned experiments & 49 \\
Result files & 793 \\
Python lines (all) & 356\,167 \\
Lean files / lines / theorems & 10 / 1\,672 / 57 \\
Manuscript & 44 pp., 51 refs \\
Ledger scripts (builders / validators) & 341 (146 / 183) \\
Ledger records (JSON / Markdown) & 171 / 186 \\
Ledger checks (top-level / package) & 4\,896 / 162\,327 \\
Manuscript validator checks & 359 \\
Memory files & 11 \\
\bottomrule
\end{tabular}
\caption{What the program produced (repository state at the end of the final phase).}
\label{tab:artefacts}
\end{table}

The frozen ledger's two validation chains take about twenty minutes, the builder sequence that
refreshes their snapshots ten, the new manuscript validator two seconds, and the Lean check under a
minute once the Mathlib cache is present. The longest single computation was the $h=10^{-4}$
double-pendulum reference over $T=3$ (1.6~h), computed once and cached.
